%% ============================================================
%%  AI Inference Engineering: Building Fast, Scalable, and
%%  Production-Ready AI Systems
%%  Chapter 1: Introduction to AI Inference Engineering
%%  Author: AI Inference Engineering Book Series
%% ============================================================

\documentclass[12pt, twoside, openright]{book}

%% -- Core encoding & fonts ----------------------------------
\usepackage[T1]{fontenc}
\usepackage[utf8]{inputenc}
\usepackage{mathpazo}   % Palatino text + math fonts (widely available)
\usepackage{microtype}

%% -- Page geometry ------------------------------------------
\usepackage[
  paperwidth=7.5in, paperheight=9.25in,
  top=1.0in, bottom=1.1in,
  inner=1.0in, outer=0.85in,
  headheight=15pt
]{geometry}

%% -- Colour palette -----------------------------------------
\usepackage[dvipsnames, svgnames, table]{xcolor}

\definecolor{chaptercolor}{RGB}{15, 82, 148}     % deep blue
\definecolor{accentorange}{RGB}{220, 100, 20}    % warm orange
\definecolor{accentgreen}{RGB}{40, 140, 80}      % emerald
\definecolor{lightbluebg}{RGB}{235, 244, 255}    % pale blue
\definecolor{codebg}{RGB}{248, 248, 252}         % near-white
\definecolor{coderule}{RGB}{180, 180, 210}       % muted purple-grey
\definecolor{keywordcolor}{RGB}{0, 100, 180}
\definecolor{commentcolor}{RGB}{100, 140, 80}
\definecolor{stringcolor}{RGB}{180, 60, 30}
\definecolor{noteboxbg}{RGB}{255, 252, 230}
\definecolor{noteboxrule}{RGB}{210, 180, 30}
\definecolor{warningbg}{RGB}{255, 240, 240}
\definecolor{warningrule}{RGB}{210, 60, 60}
\definecolor{tipbg}{RGB}{235, 255, 235}
\definecolor{tiprule}{RGB}{40, 160, 80}
\definecolor{defbg}{RGB}{240, 245, 255}
\definecolor{defrule}{RGB}{60, 100, 200}
\definecolor{tableheadbg}{RGB}{15, 82, 148}
\definecolor{tablerowodd}{RGB}{245, 249, 255}
\definecolor{tableroweven}{RGB}{255, 255, 255}

%% -- Mathematics --------------------------------------------
\usepackage{amsmath, amssymb, amsthm}

%% -- Graphics & TikZ ----------------------------------------
\usepackage{graphicx}
\usepackage{tikz}
\usepackage{pgfplots}
\pgfplotsset{compat=1.18}
\usetikzlibrary{
  arrows.meta, shapes.geometric, shapes.misc,
  positioning, fit, backgrounds, decorations.pathreplacing,
  decorations.markings, patterns, calc, shadows,
  matrix, mindmap, trees, chains
}

%% -- Tables -------------------------------------------------
\usepackage{longtable, booktabs, array, tabularx, multirow}
\usepackage{colortbl}

%% -- Code listings ------------------------------------------
\usepackage{listings}

\lstdefinestyle{pythonstyle}{
  language=Python,
  backgroundcolor=\color{codebg},
  basicstyle=\ttfamily\small\linespread{1.1}\selectfont,
  keywordstyle=\color{keywordcolor}\bfseries,
  commentstyle=\color{commentcolor}\itshape,
  stringstyle=\color{stringcolor},
  numberstyle=\tiny\color{gray},
  numbers=left,
  stepnumber=1,
  numbersep=10pt,
  frame=single,
  framerule=0.5pt,
  rulecolor=\color{coderule},
  breaklines=true,
  breakatwhitespace=true,
  tabsize=4,
  showstringspaces=false,
  captionpos=b,
  xleftmargin=18pt,
  xrightmargin=4pt,
  aboveskip=10pt,
  belowskip=6pt,
}

\lstdefinestyle{bashstyle}{
  language=bash,
  backgroundcolor=\color{black!90},
  basicstyle=\ttfamily\small\color{white},
  keywordstyle=\color{yellow!80}\bfseries,
  commentstyle=\color{green!60}\itshape,
  stringstyle=\color{orange!80},
  numbers=none,
  frame=single,
  framerule=0pt,
  breaklines=true,
  tabsize=2,
  showstringspaces=false,
  captionpos=b,
  xleftmargin=6pt,
  xrightmargin=6pt,
  aboveskip=10pt,
  belowskip=6pt,
}

\lstdefinestyle{yamlstyle}{
  language={},
  backgroundcolor=\color{codebg},
  basicstyle=\ttfamily\small,
  keywordstyle=\color{keywordcolor}\bfseries,
  commentstyle=\color{commentcolor}\itshape,
  stringstyle=\color{stringcolor},
  numbers=left,
  stepnumber=1,
  numbersep=10pt,
  frame=single,
  framerule=0.5pt,
  rulecolor=\color{coderule},
  breaklines=true,
  tabsize=2,
  showstringspaces=false,
  captionpos=b,
  xleftmargin=18pt,
  aboveskip=10pt,
  belowskip=6pt,
}

\lstset{style=pythonstyle}

%% -- Fancy boxes (tcolorbox) --------------------------------
\usepackage[most]{tcolorbox}
\tcbuselibrary{listings, breakable, skins}

%% Note box
\newtcolorbox{notebox}[1][Note]{
  enhanced, breakable,
  colback=noteboxbg, colframe=noteboxrule,
  fonttitle=\bfseries\small, title=#1,
  left=6pt, right=6pt, top=4pt, bottom=4pt,
  arc=3pt, outer arc=3pt,
  drop shadow={black!20!white, opacity=0.4},
}

%% Warning box
\newtcolorbox{warningbox}[1][Warning]{
  enhanced, breakable,
  colback=warningbg, colframe=warningrule,
  fonttitle=\bfseries\small, title=#1,
  left=6pt, right=6pt, top=4pt, bottom=4pt,
  arc=3pt,
}

%% Tip / Best-practice box
\newtcolorbox{tipbox}[1][Tip]{
  enhanced, breakable,
  colback=tipbg, colframe=tiprule,
  fonttitle=\bfseries\small, title=#1,
  left=6pt, right=6pt, top=4pt, bottom=4pt,
  arc=3pt,
}

%% Definition box
\newtcolorbox{defbox}[1][Definition]{
  enhanced, breakable,
  colback=defbg, colframe=defrule,
  fonttitle=\bfseries\small\color{white},
  coltitle=white,
  attach boxed title to top left={yshift=-2mm, xshift=6pt},
  boxed title style={colback=defrule, arc=2pt},
  title=#1,
  left=6pt, right=6pt, top=8pt, bottom=4pt,
  arc=3pt,
}

%% Analogy / Intuition box
\newtcolorbox{analogybox}[1][Intuition]{
  enhanced, breakable,
  colback=lightbluebg, colframe=chaptercolor,
  fonttitle=\bfseries\small\color{white},
  coltitle=white,
  attach boxed title to top left={yshift=-2mm, xshift=6pt},
  boxed title style={colback=chaptercolor, arc=2pt},
  title=$\star$~#1,
  left=6pt, right=6pt, top=8pt, bottom=4pt,
  arc=3pt,
}

%% Key Insight box
\newtcolorbox{insightbox}[1][Key Insight]{
  enhanced, breakable,
  colback=orange!8, colframe=accentorange,
  fonttitle=\bfseries\small\color{white},
  coltitle=white,
  attach boxed title to top left={yshift=-2mm, xshift=6pt},
  boxed title style={colback=accentorange, arc=2pt},
  title=#1,
  left=6pt, right=6pt, top=8pt, bottom=4pt,
  arc=3pt,
}

%% -- Header / Footer ----------------------------------------
\usepackage{fancyhdr}
\pagestyle{fancy}
\fancyhf{}
\fancyhead[LE]{\small\color{gray}\textit{Chapter 1 $\cdot$ Introduction to AI Inference Engineering}}
\fancyhead[RO]{\small\color{gray}\textit{AI Inference Engineering}}
\fancyfoot[C]{\small\color{gray}\thepage}
\renewcommand{\headrulewidth}{0.4pt}

%% -- Chapter/Section Titles ---------------------------------
\usepackage{titlesec}

\titleformat{\chapter}[display]
  {\normalfont\huge\bfseries\color{chaptercolor}}
  {\normalfont\large\color{chaptercolor!60}
    \MakeUppercase{\chaptertitlename}\ \thechapter}
  {16pt}
  {\Huge\bfseries}
  [\vspace{4pt}\color{chaptercolor}\rule{\textwidth}{2pt}\vspace{2pt}
   \color{chaptercolor!40}\rule{\textwidth}{0.5pt}]

\titleformat{\section}
  {\normalfont\Large\bfseries\color{chaptercolor}}
  {\thesection}{1em}{}

\titleformat{\subsection}
  {\normalfont\large\bfseries\color{chaptercolor!80}}
  {\thesubsection}{1em}{}

\titleformat{\subsubsection}
  {\normalfont\normalsize\bfseries\color{accentorange}}
  {\thesubsubsection}{1em}{}

%% -- Lists --------------------------------------------------
\usepackage{enumitem}
\setlist[itemize]{leftmargin=*, itemsep=3pt, topsep=4pt}
\setlist[enumerate]{leftmargin=*, itemsep=3pt, topsep=4pt}

%% -- Hyperref (load last) -----------------------------------
\usepackage{hyperref}
\hypersetup{
  colorlinks=true,
  linkcolor=chaptercolor,
  citecolor=accentgreen,
  urlcolor=accentorange,
  pdftitle={AI Inference Engineering -- Chapter 1},
  pdfauthor={AI Inference Engineering Book Series},
  pdfsubject={AI Inference Engineering},
  bookmarksnumbered=true,
  bookmarksopen=true,
}

%% -- Utility macros -----------------------------------------
\newcommand{\term}[1]{\textbf{\color{chaptercolor}#1}}
\newcommand{\code}[1]{\texttt{\small\color{keywordcolor}#1}}
\newcommand{\metric}[1]{\textbf{\color{accentorange}#1}}
\newcommand{\chapquote}[2]{%
  \begin{center}
    \begin{minipage}{0.82\textwidth}
      \itshape\large ``#1''
      \vspace{4pt}\\
      \normalsize\normalfont\hfill --- #2
    \end{minipage}
  \end{center}
  \vspace{8pt}
}

%% -- Learning-objectives environment ------------------------
\newenvironment{learningobjectives}{%
  \begin{tcolorbox}[
    enhanced, breakable,
    colback=chaptercolor!6, colframe=chaptercolor,
    title={\bfseries\color{white} Learning Objectives},
    fonttitle=\bfseries,
    coltitle=white,
    colbacktitle=chaptercolor,
    arc=4pt,
    left=8pt, right=8pt, top=6pt, bottom=6pt,
  ]
  \begin{enumerate}[label=\textcolor{chaptercolor}{\bfseries\arabic*.}, leftmargin=*]
}{%
  \end{enumerate}
  \end{tcolorbox}
}

%% -- Exercise environment -----------------------------------
\newcounter{exercise}[chapter]
\newenvironment{exercise}[1][]{%
  \refstepcounter{exercise}
  \begin{tcolorbox}[
    enhanced,
    colback=accentgreen!5, colframe=accentgreen,
    title={\bfseries\color{white} Exercise \theexercise\ifx&#1&\else: #1\fi},
    coltitle=white, colbacktitle=accentgreen,
    arc=3pt, left=6pt, right=6pt, top=6pt, bottom=6pt,
  ]
}{%
  \end{tcolorbox}
}

%% -- Further Reading environment ----------------------------
\newenvironment{furtherreading}{%
  \begin{tcolorbox}[
    enhanced, breakable,
    colback=gray!5, colframe=gray!50,
    title={\bfseries Further Reading \& References},
    coltitle=black, colbacktitle=gray!20,
    arc=3pt, left=6pt, right=6pt,
  ]
}{%
  \end{tcolorbox}
}

%% -- Summary box --------------------------------------------
\newenvironment{chapsummary}{%
  \begin{tcolorbox}[
    enhanced, breakable,
    colback=chaptercolor!5, colframe=chaptercolor,
    title={\bfseries\color{white} Chapter Summary},
    coltitle=white, colbacktitle=chaptercolor,
    arc=4pt, left=8pt, right=8pt, top=6pt, bottom=6pt,
  ]
}{%
  \end{tcolorbox}
}

%% -----------------------------------------------------------
%%  DOCUMENT START
%% -----------------------------------------------------------
\begin{document}

\frontmatter

%% -- Cover Page ---------------------------------------------
\begin{titlepage}
\pagecolor{chaptercolor}
\color{white}
\centering
\vspace*{3cm}

{\fontsize{13}{16}\selectfont\MakeUppercase{AI Inference Engineering}}

\vspace{0.6cm}
\color{white!70}\rule{0.6\textwidth}{0.5pt}\color{white}
\vspace{0.6cm}

{\fontsize{22}{28}\selectfont\bfseries Building Fast, Scalable, and\\[4pt]
Production-Ready AI Systems}

\vspace{2.2cm}

{\fontsize{16}{20}\selectfont\color{orange!80}
\textbf{Chapter 1}}

\vspace{0.4cm}

{\fontsize{24}{30}\selectfont\bfseries
Introduction to\\[6pt] AI Inference Engineering}

\vspace{3.5cm}

\begin{tikzpicture}[scale=0.9]
  % GPU chip diagram
  \foreach \x in {0,1,2,3} {
    \foreach \y in {0,1,2,3} {
      \fill[white!20] (\x*1.1, \y*1.1) rectangle (\x*1.1+0.9, \y*1.1+0.9);
      \fill[orange!60] (\x*1.1+0.15, \y*1.1+0.15)
                       rectangle (\x*1.1+0.75, \y*1.1+0.75);
    }
  }
  \draw[white, very thick] (-0.15,-0.15) rectangle (4.75,4.75);
  \node[white, font=\footnotesize\bfseries] at (2.3, -0.55) {GPU Compute Grid};
  % Arrows indicating data flow
  \draw[->, white!60, very thick] (5.2, 2.3) -- (4.9, 2.3);
  \node[white!80, font=\tiny, right] at (5.25, 2.45) {Input};
  \draw[->, white!60, very thick] (-0.3, 2.3) -- (-0.6, 2.3);
  \node[white!80, font=\tiny, left] at (-0.65, 2.45) {Output};
\end{tikzpicture}

\vspace{2.5cm}
\color{white!60}
{\small\itshape First Edition $\cdot$ 2025}

\end{titlepage}
\nopagecolor
\color{black}

\mainmatter
\setcounter{chapter}{0}

%% -----------------------------------------------------------
\chapter{Introduction to AI Inference Engineering}
%% -----------------------------------------------------------

\chapquote{%
  Any sufficiently advanced technology is indistinguishable from magic.
  Our job as engineers is to make that magic fast, cheap, and reliable.%
}{Arthur C. Clarke (adapted)}

\vspace{0.5cm}

\begin{learningobjectives}
  \item Understand the precise meaning of \textbf{inference} and how it
        differs from model training.
  \item Explain why inference engineering has become one of the most
        important disciplines in modern AI infrastructure.
  \item Describe the full lifecycle of a request through a production
        inference system --- from API call to token response.
  \item Understand the key performance metrics used to evaluate inference
        systems: \textbf{latency}, \textbf{throughput}, and
        \textbf{cost per token}.
  \item Identify the main hardware components (GPU, CPU, memory, network)
        and their roles in inference.
  \item Trace the historical evolution from research models to today's
        hyper-scale inference platforms.
  \item Recognise the real-world systems --- ChatGPT, Gemini, Claude,
        Midjourney --- that rely on the engineering concepts taught in
        this book.
\end{learningobjectives}

%% -----------------------------------------------------------
\section{Why This Book Exists}
%% -----------------------------------------------------------

Every time you send a message to ChatGPT, ask Gemini to summarise a
document, or watch Midjourney paint an image from your words, something
remarkable is happening behind the scenes.

Thousands of specialised computers --- each containing powerful graphics
processors, hundreds of gigabytes of memory, and high-speed networking ---
are working together at millisecond timescales to run an AI model and
deliver your result.

That entire operation is called \term{AI inference}. The engineering
discipline that designs, builds, optimises, and scales these systems is
called \term{AI inference engineering}.

This book will teach you that discipline --- from the ground up.

\begin{notebox}[Who This Book Is For]
This book is written for:
\begin{itemize}
  \item Software engineers who want to move into AI infrastructure
  \item ML engineers who want to understand what happens after model
        training
  \item Platform engineers building internal AI serving systems
  \item Technical leads making architecture decisions about AI deployments
  \item Anyone curious about how large AI systems actually work at scale
\end{itemize}
You do \textbf{not} need a PhD. You do need curiosity and a willingness
to think carefully about systems.
\end{notebox}

%% -----------------------------------------------------------
\section{What Is Inference? A First-Principles Explanation}
%% -----------------------------------------------------------

\subsection{Start with the Simplest Idea}

Let's start with something familiar.

Imagine you study very hard for a maths exam.
You spend weeks reading textbooks, working through problems,
and learning formulas.
That studying phase --- where you absorb knowledge --- is the
\textit{learning} phase.

Then exam day arrives.
Someone gives you a new problem you have never seen before.
You apply everything you have learned to produce an answer.
That answering phase is \textit{inference}.

\begin{analogybox}[The Exam Analogy]
\textbf{Training} = studying for months to learn how to solve problems\\[4pt]
\textbf{Inference} = sitting the exam and solving a new problem using
what you learned\\[6pt]
The student does not re-read all their textbooks during the exam.
They use their \textit{trained knowledge} --- stored in their brain ---
to answer quickly.\\[4pt]
AI systems work the same way.
\end{analogybox}

\subsection{Translating the Analogy to AI}

An AI language model like GPT-4 or Claude is trained on hundreds of
billions of text documents.
During training, the model adjusts billions of numerical parameters
--- think of these as the ``knowledge'' stored in the model ---
to learn patterns in language.

This training process takes \textit{months}, thousands of GPUs,
and costs tens of millions of dollars.

Once training is complete, the model's parameters are frozen.
The model is then \textit{deployed} --- made available to users.

When you send a message to ChatGPT:

\begin{enumerate}
  \item Your message arrives at a server.
  \item The server loads the trained model parameters.
  \item The model performs a mathematical computation on your message
        using those parameters.
  \item The model outputs a response --- word by word, or technically,
        \textit{token by token}.
  \item The response is sent back to you.
\end{enumerate}

That entire process --- steps 1 through 5 --- is called a single
\term{inference request}.

\begin{defbox}[AI Inference]
\textbf{AI Inference} is the process of running a trained AI model
on new input data to produce a prediction or output.\\[4pt]
The model parameters (weights) do not change during inference.
The model simply \textit{applies} its learned knowledge to the input.
\end{defbox}

\subsection{The Key Mathematical Operation}

If you want to understand inference at its heart, here is the core
mathematical reality:

An AI language model is, at its most fundamental level, a very large
mathematical function:

\[
  \mathbf{output} = f(\mathbf{input},\ \boldsymbol{\theta})
\]

where:
\begin{itemize}
  \item $\mathbf{input}$ is the text you provided (encoded as numbers)
  \item $\boldsymbol{\theta}$ (theta) represents the model's
        \textit{billions of parameters} --- its learned knowledge
  \item $f$ is the series of mathematical operations the model performs
        (matrix multiplications, attention computations, activations)
  \item $\mathbf{output}$ is the probability distribution over the next
        token to generate
\end{itemize}

During inference, $\boldsymbol{\theta}$ is \textbf{fixed}.
Only $\mathbf{input}$ varies.
The model's job is simply to evaluate this function quickly and correctly.

\begin{insightbox}[The Core Challenge]
GPT-4 has approximately \textbf{1.76 trillion parameters}.
Each parameter is a floating-point number (typically 16 bits = 2 bytes).
That alone is \textbf{3.5 terabytes of data} --- just to store the model.\\[4pt]
Loading, managing, and computing with that much data,
at high speed, for millions of users simultaneously,
is the fundamental engineering challenge of AI inference.
\end{insightbox}

%% -----------------------------------------------------------
\section{Training vs.\ Inference: Understanding the Difference}
%% -----------------------------------------------------------

This is one of the most important distinctions in AI engineering.
Training and inference look superficially similar --- both run the
model on data --- but they are \textit{fundamentally different
engineering problems}.

\subsection{What Happens During Training?}

During training, the goal is to \textit{improve} the model's parameters.

The process works like this:
\begin{enumerate}
  \item Feed a batch of training examples into the model.
  \item The model makes predictions (a \textit{forward pass}).
  \item Compare the predictions to the correct answers.
        Compute a \textit{loss} (how wrong was the model?).
  \item Use calculus (\textit{backpropagation}) to compute how each
        parameter contributed to the error.
  \item Adjust the parameters slightly in the direction that reduces
        the error (\textit{gradient descent}).
  \item Repeat billions of times.
\end{enumerate}

Step 4 --- backpropagation --- requires storing intermediate computations
from the forward pass so gradients can flow backwards.
This is called the \term{activation memory} or \term{gradient memory}
and it is \textit{enormous} --- often 5--10x the size of the model itself.

\subsection{What Happens During Inference?}

During inference:
\begin{enumerate}
  \item Feed a user's input into the model.
  \item The model makes a prediction (forward pass only).
  \item Return the prediction.
\end{enumerate}

That's it. No backpropagation. No gradient storage. No parameter updates.

\begin{figure}[htbp]
\centering
\begin{tikzpicture}[
  node distance=0.7cm and 1.6cm,
  block/.style={
    rectangle, rounded corners=4pt, draw,
    minimum width=2.8cm, minimum height=0.8cm,
    font=\small, align=center
  },
  trainblock/.style={block, fill=orange!20, draw=orange!70},
  inferblock/.style={block, fill=blue!15, draw=blue!60},
  arrow/.style={->, >=Stealth, thick},
  trainarrow/.style={arrow, color=orange!80},
  inferarrow/.style={arrow, color=blue!70},
]

%% TRAINING column
\node[trainblock] (tdata)  {Training Data};
\node[trainblock, below=of tdata] (tforward) {Forward Pass\\{\tiny compute predictions}};
\node[trainblock, below=of tforward] (tloss)   {Loss Computation\\{\tiny how wrong?}};
\node[trainblock, below=of tloss] (tback)    {Backpropagation\\{\tiny compute gradients}};
\node[trainblock, below=of tback] (tupdate)  {Parameter Update\\{\tiny improve model}};

\draw[trainarrow] (tdata) -- (tforward);
\draw[trainarrow] (tforward) -- (tloss);
\draw[trainarrow] (tloss) -- (tback);
\draw[trainarrow] (tback) -- (tupdate);
\draw[trainarrow] (tupdate.east) to[out=0, in=0, looseness=2.5]
                  node[right, font=\tiny\color{orange!80}]{repeat / billions of times}
                  (tdata.east);

\node[above=0.2cm of tdata, font=\bfseries\color{orange!80}] {TRAINING};
\node[below=0.2cm of tupdate, font=\small\color{orange!60},
      text width=3.5cm, align=center]
  {Parameters \textbf{change}\\weeks on 1000s of GPUs};

%% INFERENCE column
\node[inferblock, right=3.8cm of tdata] (iinput) {User Input};
\node[inferblock, below=of iinput] (iforward) {Forward Pass\\{\tiny apply model}};
\node[inferblock, below=of iforward] (ioutput)  {Output Tokens\\{\tiny prediction}};
\node[inferblock, below=of ioutput] (ireturn) {Return Response\\{\tiny to user}};

\draw[inferarrow] (iinput) -- (iforward);
\draw[inferarrow] (iforward) -- (ioutput);
\draw[inferarrow] (ioutput) -- (ireturn);

\node[above=0.2cm of iinput, font=\bfseries\color{blue!70}] {INFERENCE};
\node[below=0.2cm of ireturn, font=\small\color{blue!60},
      text width=3.5cm, align=center]
  {Parameters \textbf{fixed}\\milliseconds per request};

%% Vertical separator
\draw[dashed, gray!50, thick]
  ($(tdata.east)!0.5!(iinput.west)+(0,0.8)$) --
  ($(tupdate.east)!0.5!(ireturn.west)-(0,0.8)$);

%% Memory annotations
\draw[<->, red!60, dashed]
  (tback.west) -- ++(-0.7,0)
  node[left, red!60, font=\tiny, text width=1.4cm, align=center]
    {Stores gradients\\5--10$\times$ model size};

\draw[<->, accentgreen!70, dashed]
  (iforward.east) -- ++(0.7,0)
  node[right, accentgreen!70, font=\tiny, text width=1.4cm, align=center]
    {No gradients\\stored};

\end{tikzpicture}
\caption{Training vs.\ Inference. Training requires storing large
intermediate activations for backpropagation. Inference only needs
a single forward pass, making memory requirements dramatically lower.}
\label{fig:train-vs-infer}
\end{figure}

\subsection{Why This Distinction Matters Enormously}

The difference between training and inference isn't just academic ---
it completely changes the engineering requirements:

\vspace{6pt}
\rowcolors{2}{tablerowodd}{tableroweven}
\begin{longtable}{>{\bfseries}p{3cm} p{5cm} p{5cm}}
\rowcolor{tableheadbg}
\textcolor{white}{\textbf{Dimension}} &
\textcolor{white}{\textbf{Training}} &
\textcolor{white}{\textbf{Inference}} \\
\toprule
\endfirsthead
\rowcolor{tableheadbg}
\textcolor{white}{\textbf{Dimension}} &
\textcolor{white}{\textbf{Training}} &
\textcolor{white}{\textbf{Inference}} \\
\toprule
\endhead
Goal & Improve model weights & Apply model to new inputs \\
Frequency & Once (or periodic) & Millions of times per day \\
Latency need & Hours/days acceptable & Milliseconds required \\
Memory use & Very high (gradients) & Lower (weights only) \\
Compute pattern & Batch, parallel & Often sequential (autoregressive) \\
Cost driver & GPU-hours of compute & Cost per token at scale \\
Batching & Large fixed batches & Dynamic, variable batch sizes \\
Parameter state & Changing constantly & Completely frozen \\
Primary metric & Loss convergence & Latency / throughput \\
Engineering focus & Distributed training & Serving optimisation \\
\bottomrule
\caption{Key differences between training and inference workloads.}
\label{tab:train-vs-infer}
\end{longtable}

\begin{notebox}[The Business Reality]
OpenAI reportedly spends approximately \$700,000 per day running
ChatGPT inference. Google runs billions of AI inference calls daily
across its products. Amazon, Microsoft, and Meta each operate
inference clusters of tens of thousands of GPUs.\\[4pt]
Training is a one-time (or infrequent) cost.
\textbf{Inference is an ongoing, relentless operational cost.}
This is why inference engineering has become a separate discipline.
\end{notebox}

%% -----------------------------------------------------------
\section{Why Inference Engineering Matters}
%% -----------------------------------------------------------

You might wonder: ``If inference is just running the model,
how complex can it really be?''

The answer is: extraordinarily complex. Let us walk through why.

\subsection{The Scale Problem}

Consider ChatGPT during peak usage:
\begin{itemize}
  \item Approximately \textbf{10 million daily active users}
  \item Each user sends multiple messages per session
  \item Each message might generate 200--500 output tokens
  \item Each token requires evaluating a model with hundreds of billions
        of parameters
  \item Total: potentially \textbf{billions of model evaluations per day}
\end{itemize}

Now add this constraint: users expect responses within \textbf{1--2 seconds}.

This is roughly equivalent to a restaurant that must:
\begin{itemize}
  \item Serve \textit{millions of customers simultaneously}
  \item Each meal is a custom five-course dinner (not fast food)
  \item Every meal must be ready within 2 minutes of ordering
  \item The kitchen has a fixed, enormous cost to operate
\end{itemize}

No traditional software system faces this combination of constraints.
Inference engineering exists to solve them.

\subsection{The Three Core Tensions}

Every inference engineering decision involves balancing three competing
objectives:

\begin{figure}[htbp]
\centering
\begin{tikzpicture}[scale=1.0]
  % Triangle
  \coordinate (A) at (0, 0);
  \coordinate (B) at (6, 0);
  \coordinate (C) at (3, 5.2);

  \fill[chaptercolor!10] (A) -- (B) -- (C) -- cycle;
  \draw[chaptercolor, very thick] (A) -- (B) -- (C) -- cycle;

  % Labels at corners
  \node[below left, font=\bfseries\large\color{chaptercolor},
        align=center] at (A)
    {Low\\Latency};
  \node[below right, font=\bfseries\large\color{chaptercolor},
        align=center] at (B)
    {High\\Throughput};
  \node[above, font=\bfseries\large\color{chaptercolor},
        align=center] at (C)
    {Low\\Cost};

  % Central annotation
  \node[font=\small\itshape, text width=3cm, align=center,
        color=accentorange] at (3, 1.8)
    {You can only fully\\optimise \textbf{two} of these\\at once};

  % Small icons
  \node[font=\Large] at (0.2, 0.5) {\textbf{L}};
  \node[font=\Large] at (5.8, 0.5) {\textbf{T}};
  \node[font=\Large] at (3, 4.6)   {\textbf{C}};

  % Edge annotations
  \node[font=\small\color{gray}, rotate=41, above] at (1.5, 2.6)
    {Batching helps};
  \node[font=\small\color{gray}, rotate=-41, above] at (4.5, 2.6)
    {Quantisation helps};
  \node[font=\small\color{gray}, below] at (3, -0.3)
    {Caching helps};

\end{tikzpicture}
\caption{The Inference Engineering Trilemma. Every system design must
explicitly choose which objectives to prioritise.}
\label{fig:trilemma}
\end{figure}

\begin{itemize}
  \item \term{Latency}: How long does a single user wait for their
        response? (measured in milliseconds)
  \item \term{Throughput}: How many total requests can the system
        handle per second across all users?
  \item \term{Cost}: How much does it cost to serve each request?
        (measured in dollars per million tokens, often written
        \$/MTok)
\end{itemize}

These three objectives \textit{pull against each other}.
Making one better often makes another worse.

\begin{analogybox}[The Bus vs.\ Taxi Analogy]
\textbf{A taxi} picks you up immediately (low latency) but can only
carry one passenger (low throughput) and is expensive per person
(high cost).\\[4pt]
\textbf{A bus} waits until it is full (higher latency) but carries
many passengers at once (high throughput) and is cheap per person
(low cost).\\[4pt]
Inference systems face exactly this tradeoff.
\textit{Batching} is the bus strategy.
\textit{Streaming} is the taxi strategy.
The best inference engines choose intelligently based on load.
\end{analogybox}

\subsection{Performance Numbers That Matter in Production}

Here are the performance targets that production AI companies actually
care about:

\vspace{6pt}
\rowcolors{2}{tablerowodd}{tableroweven}
\begin{longtable}{>{\bfseries}p{4cm} p{4cm} p{4cm}}
\rowcolor{tableheadbg}
\textcolor{white}{\textbf{Metric}} &
\textcolor{white}{\textbf{Target (Consumer)}} &
\textcolor{white}{\textbf{Target (Enterprise)}} \\
\toprule
\endfirsthead
\rowcolor{tableheadbg}
\textcolor{white}{\textbf{Metric}} &
\textcolor{white}{\textbf{Target (Consumer)}} &
\textcolor{white}{\textbf{Target (Enterprise)}} \\
\toprule
\endhead
Time to first token (TTFT) & $<$ 500ms & $<$ 200ms \\
Tokens per second (TPS) & 20--50 tok/s & 50--200 tok/s \\
P99 latency & $<$ 5s & $<$ 2s \\
Cost per million tokens & \$0.50--\$5.00 & \$0.10--\$1.00 \\
GPU utilisation & $>$ 60\% & $>$ 80\% \\
Requests per GPU per hour & 100--500 & 500--5000 \\
\bottomrule
\caption{Typical production performance targets for LLM inference systems.}
\label{tab:perf-targets}
\end{longtable}

\begin{notebox}[What TTFT Means]
\term{Time to First Token (TTFT)} is the time from when you press
``send'' to when you see the \textit{first word} of the response appear.
It is arguably the most important user-facing latency metric because
it determines whether the system feels ``responsive'' or ``slow.''\\[4pt]
A system generating 50 tokens per second but with a 3-second TTFT
will feel slow, even if the total generation time is short.
\end{notebox}

%% -----------------------------------------------------------
\section{The Anatomy of an Inference System}
%% -----------------------------------------------------------

Before we go deeper, let us trace exactly what happens when you send
a message to a production AI system.
Understanding this flow is essential --- every optimisation technique
we will learn in later chapters targets a specific step in this journey.

\subsection{The Complete Request Lifecycle}

\begin{figure}[htbp]
\centering
\begin{tikzpicture}[
  node distance=0.5cm and 0.2cm,
  userbox/.style={
    rectangle, rounded corners=6pt, draw=accentorange!80,
    fill=accentorange!10, minimum width=2.0cm, minimum height=0.7cm,
    font=\small\bfseries, align=center
  },
  sysbox/.style={
    rectangle, rounded corners=4pt, draw=chaptercolor!70,
    fill=chaptercolor!8, minimum width=2.5cm, minimum height=0.7cm,
    font=\small, align=center
  },
  gpubox/.style={
    rectangle, rounded corners=4pt, draw=accentgreen!80,
    fill=accentgreen!10, minimum width=2.5cm, minimum height=0.7cm,
    font=\small, align=center
  },
  arrow/.style={->, >=Stealth, thick, color=gray!70},
  label/.style={font=\tiny\color{gray}, midway, above},
]

% Row 1: User
\node[userbox] (user) {User\\(Browser/App)};

% Row 2: Edge
\node[sysbox, right=1.6cm of user] (cdn) {CDN /\\Load Balancer};
\node[sysbox, right=1.0cm of cdn] (gateway) {API\\Gateway};
\node[sysbox, right=1.0cm of gateway] (auth) {Auth \&\\Rate Limit};

% Row 3: Inference engine
\node[sysbox, below=1.6cm of cdn] (scheduler) {Request\\Scheduler};
\node[sysbox, right=1.0cm of scheduler] (tokenizer) {Tokeniser\\Preprocessing};
\node[gpubox, right=1.0cm of tokenizer] (prefill) {Prefill\\(GPU)};
\node[gpubox, right=1.0cm of prefill] (decode) {Decode\\(GPU loop)};

% Row 4: Storage
\node[sysbox, below=1.6cm of tokenizer] (kvcache) {KV Cache\\(GPU Memory)};
\node[sysbox, right=1.0cm of kvcache] (weights) {Model Weights\\(GPU Memory)};
\node[sysbox, right=1.0cm of weights] (detok) {Detokeniser\\Postprocessing};

% Arrows - top row
\draw[arrow] (user) -- (cdn) node[label] {HTTPS};
\draw[arrow] (cdn) -- (gateway) node[label] {HTTP/2};
\draw[arrow] (gateway) -- (auth) node[label] {};

% Down to scheduler
\draw[arrow] (auth) -- ++(0,-0.9) -| (scheduler);

% Middle row
\draw[arrow] (scheduler) -- (tokenizer) node[label] {};
\draw[arrow] (tokenizer) -- (prefill) node[label] {};
\draw[arrow] (prefill) -- (decode) node[label] {KV cache};

% Down to storage
\draw[arrow] (prefill.south) -- ++(0,-0.4) -| (kvcache);
\draw[arrow] (kvcache.north) -- ++(0,0.4) -| (prefill);

% Weights to GPUs
\draw[arrow] (weights.north) -- ++(0,0.8) -| (prefill.south);

% Decode to detokeniser
\draw[arrow] (decode.south) -- ++(0,-0.8) -- (detok.north);

% Back to user
\draw[arrow, color=accentorange!60] (detok.west) -- ++(-8.5,0) node[midway, above, font=\tiny\color{accentorange}] {streaming tokens back to user} |- (user.south);

% Timing annotations
\node[font=\tiny, color=gray, text width=1.4cm, align=center,
      below=0.1cm of scheduler]
  {$<$1ms routing};
\node[font=\tiny, color=gray, text width=1.4cm, align=center,
      below=0.1cm of prefill]
  {10--200ms\\prefill};
\node[font=\tiny, color=gray, text width=1.4cm, align=center,
      below=0.1cm of decode]
  {20ms/token\\decode};

\end{tikzpicture}
\caption{Anatomy of a production inference request. Each box represents
a system component that we will study in detail in later chapters.}
\label{fig:request-lifecycle}
\end{figure}

Let's walk through each stage:

\subsubsection{Stage 1: Network Ingress (0--5 ms)}
Your request travels over the internet to the nearest \term{CDN edge
node} (Content Delivery Network). The request is routed to the closest
data centre to reduce network round-trip time.

\subsubsection{Stage 2: API Gateway and Authentication (1--10 ms)}
The API gateway validates your API key, checks your rate limits,
and decides whether to accept or queue your request.
This protects the system from being overwhelmed.

\subsubsection{Stage 3: Request Scheduling ($<$1 ms)}
The \term{request scheduler} --- one of the most important components
in a modern inference engine --- decides:
\begin{itemize}
  \item Which GPU cluster will handle this request?
  \item Should this request be batched with other pending requests?
  \item How much priority should this request receive?
\end{itemize}

\subsubsection{Stage 4: Tokenisation (0.5--2 ms)}
Your text is converted into \term{tokens} --- numerical IDs from the
model's vocabulary. The word ``inference'' might become token 23847.
``engineering'' might become tokens 45231 + 9182 (split into sub-words).
This tokenised input is what the GPU will actually process.

\subsubsection{Stage 5: Prefill Phase (10--500 ms)}
This is where real GPU computation begins.
The model processes your \textit{entire input prompt} in parallel to
build an internal representation of what you have said.

Think of it as the model ``reading and comprehending'' your entire
message at once before it starts responding.

This phase is compute-intensive and its duration scales with the
length of your prompt.

\subsubsection{Stage 6: Decode Phase (20--50 ms per token)}
Now the model begins generating tokens \textit{one at a time}.
Each token requires:
\begin{enumerate}
  \item Loading the model weights from GPU memory
  \item Performing a forward pass through the model
  \item Sampling the next token from the probability distribution
  \item Appending the new token to the context
  \item Repeating for the next token
\end{enumerate}

This sequential nature is what makes LLM inference fundamentally
different from other AI workloads --- and is the source of many of
the key optimisations we will study.

\subsubsection{Stage 7: Detokenisation and Return}
Each generated token ID is converted back to text and streamed back
to you --- which is why you see text appearing word by word on your
screen rather than all at once.

\begin{insightbox}[The Decode Phase Bottleneck]
The decode phase is typically the \textbf{primary bottleneck} in
LLM inference. Why? Because each token generation requires reading
the \textit{entire model} from GPU memory --- often hundreds of
gigabytes. Modern GPUs can read memory at 3--8 TB/s, but that is
still the limiting factor. This is called being
\textbf{memory-bandwidth-bound}, and we will explore it extensively
in Chapter 3.
\end{insightbox}

%% -----------------------------------------------------------
\section{The Hardware Landscape}
%% -----------------------------------------------------------

AI inference runs on specialised hardware. Understanding this hardware
--- at least conceptually --- is essential for understanding why
inference systems are designed the way they are.

\subsection{The GPU: Heart of Inference}

A \term{GPU} (Graphics Processing Unit) was originally designed to
render video game graphics. It turns out that the mathematical
operations needed for graphics --- thousands of small matrix
multiplications in parallel --- are identical to the operations
needed for AI neural network inference.

\begin{analogybox}[CPU vs.\ GPU: The Factory Analogy]
\textbf{A CPU} is like a small team of highly-skilled engineers.
There are 8--64 of them. Each one is incredibly smart, can handle
complex tasks, and can switch between very different jobs instantly.
But you only have a few of them.\\[6pt]
\textbf{A GPU} is like a factory floor with 10,000 workers.
Each individual worker is less capable --- they can only do simple
arithmetic. But having 10,000 of them working simultaneously means
they can process enormous amounts of data in parallel.\\[6pt]
Neural network inference is essentially: ``multiply these 100 billion
numbers together in the right way.'' That's a perfect job for 10,000
parallel workers.
\end{analogybox}

\subsection{Key GPU Specifications That Matter}

When you hear about an NVIDIA H100 or A100 GPU being used for AI
inference, here are the numbers that actually matter:

\vspace{6pt}
\rowcolors{2}{tablerowodd}{tableroweven}
\begin{longtable}{>{\bfseries}p{3.5cm} p{3.2cm} p{3.2cm} p{2.5cm}}
\rowcolor{tableheadbg}
\textcolor{white}{\textbf{Specification}} &
\textcolor{white}{\textbf{NVIDIA A100 (80G)}} &
\textcolor{white}{\textbf{NVIDIA H100 (80G)}} &
\textcolor{white}{\textbf{Why It Matters}} \\
\toprule
\endfirsthead
\rowcolor{tableheadbg}
\textcolor{white}{\textbf{Specification}} &
\textcolor{white}{\textbf{NVIDIA A100 (80G)}} &
\textcolor{white}{\textbf{NVIDIA H100 (80G)}} &
\textcolor{white}{\textbf{Why It Matters}} \\
\toprule
\endhead
VRAM (memory) & 80 GB & 80 GB & Model must fit in VRAM \\
Memory bandwidth & 2.0 TB/s & 3.35 TB/s & Decode speed \\
FP16 TFLOPS & 312 & 989 & Prefill speed \\
Tensor cores & 432 & 528 & Matrix multiply speed \\
NVLink bandwidth & 600 GB/s & 900 GB/s & Multi-GPU communication \\
TDP (power) & 400W & 700W & Operational cost \\
\bottomrule
\caption{Key GPU specifications relevant to AI inference.}
\label{tab:gpu-specs}
\end{longtable}

\subsection{The Memory Hierarchy}

Understanding where data lives during inference is crucial.
The memory hierarchy is a pyramid: faster memory at the top but
smaller capacity, slower memory at the bottom but larger capacity.

\begin{figure}[htbp]
\centering
\begin{tikzpicture}[scale=0.95]
  % Define pyramid layers from top to bottom
  % Each layer: y-position, half-width, colour, label, stats

  % Layer 5 (top): Registers / L1 Cache
  \fill[red!80] (0,6.5) -- (0.6,7.2) -- (-0.6,7.2) -- cycle;
  \fill[red!80] (-0.6,6.5) rectangle (0.6,7.2);
  \node[right, font=\small\bfseries] at (0.7, 6.85)
    {L1 Cache / Shared Memory};
  \node[right, font=\tiny\color{gray}] at (0.7, 6.55)
    {$\sim$256 KB/SM\ \ |\ \ $\sim$19 TB/s};

  % Layer 4: L2 Cache
  \fill[orange!70] (-1.2,5.7) rectangle (1.2,6.4);
  \node[right, font=\small\bfseries] at (1.3, 6.05)
    {L2 Cache};
  \node[right, font=\tiny\color{gray}] at (1.3, 5.75)
    {50 MB\ \ |\ \ $\sim$7 TB/s};

  % Layer 3: HBM (GPU VRAM)
  \fill[yellow!70] (-2.2,4.5) rectangle (2.2,5.6);
  \node[right, font=\small\bfseries] at (2.3, 5.2)
    {HBM (GPU VRAM)};
  \node[right, font=\tiny\color{gray}] at (2.3, 4.7)
    {80 GB\ \ |\ \ 3.35 TB/s (H100)};
  \node[font=\tiny\itshape, text width=4.3cm, align=center,
        color=chaptercolor] at (0,5.05)
    {Model weights live here\\KV cache lives here};

  % Layer 2: CPU DRAM / NVLink
  \fill[accentgreen!40] (-3.2,3.2) rectangle (3.2,4.4);
  \node[right, font=\small\bfseries] at (3.3, 3.95)
    {CPU DRAM / NVLink};
  \node[right, font=\tiny\color{gray}] at (3.3, 3.45)
    {0.5--2 TB\ \ |\ \ 50--900 GB/s};
  \node[font=\tiny\itshape, color=gray] at (0,3.8)
    {Overflow weights, inter-GPU data};

  % Layer 1 (bottom): SSD / NVMe
  \fill[blue!25] (-4.2,1.7) rectangle (4.2,3.1);
  \node[right, font=\small\bfseries] at (4.3, 2.55)
    {NVMe SSD};
  \node[right, font=\tiny\color{gray}] at (4.3, 2.05)
    {10--100 TB\ \ |\ \ 7--12 GB/s};
  \node[font=\tiny\itshape, color=gray] at (0,2.4)
    {Model checkpoint storage};

  % Layer 0: Object Storage
  \fill[gray!20] (-5.0,0.2) rectangle (5.0,1.6);
  \node[right, font=\small\bfseries] at (5.1, 1.1)
    {Object Storage (S3/GCS)};
  \node[right, font=\tiny\color{gray}] at (5.1, 0.6)
    {Unlimited\ \ |\ \ 1--10 GB/s};

  % Speed labels on left
  \node[left, font=\tiny\bfseries\color{red!70}] at (-4.3, 6.85) {FASTEST};
  \node[left, font=\tiny\bfseries\color{gray}] at (-4.3, 0.9)    {SLOWEST};
  \draw[->, gray!50, dashed] (-4.1,1.0) -- (-4.1,6.7);
  \node[left, font=\tiny\color{gray}, rotate=90] at (-4.6,3.9) {Speed};

\end{tikzpicture}
\caption{The memory hierarchy in a modern GPU inference system.
The goal of inference optimisation is to keep hot data as high up
this pyramid as possible.}
\label{fig:memory-hierarchy}
\end{figure}

The fundamental rule of inference engineering is:
\textbf{data that is accessed frequently should live as close to the
GPU compute units as possible.}

Every optimisation technique we study --- KV caching, quantisation,
weight sharing --- can be understood as trying to keep the right data
at the right level of this memory hierarchy.

%% -----------------------------------------------------------
\section{The Evolution of AI Inference Infrastructure}
%% -----------------------------------------------------------

To understand where we are today and where we are going, we need to
understand where we came from.

\subsection{Era 1: Research-Scale Models (2012--2017)}

The deep learning revolution began with AlexNet winning the ImageNet
competition in 2012. Models were relatively small (millions of parameters),
ran on single GPUs, and inference was simple: run the model, get the
prediction, done.

The engineering challenge was primarily training efficiency.
Inference was an afterthought.

\subsection{Era 2: The Rise of Pre-trained Models (2018--2020)}

BERT (2018) and GPT-2 (2019) introduced the concept of large pre-trained
transformer models fine-tuned for downstream tasks.
Models grew to billions of parameters.

Suddenly, inference required serious thought:
\begin{itemize}
  \item Models no longer fit easily on single GPUs
  \item Serving latency became a customer-facing concern
  \item Companies began building dedicated serving infrastructure
\end{itemize}

TensorFlow Serving and TorchServe emerged as early inference servers.

\subsection{Era 3: The LLM Explosion (2020--2022)}

GPT-3 (2020) with 175 billion parameters changed everything.
No single consumer GPU could hold it.
Running inference required \textit{model parallelism} across multiple GPUs.

OpenAI, Google, and DeepMind began building specialised inference
infrastructure at massive scale.
The \term{transformer} architecture became universal.
New optimisation techniques emerged:
\begin{itemize}
  \item Quantisation (reduce precision to save memory)
  \item KV caching (avoid redundant computation)
  \item Tensor parallelism (split the model across GPUs)
  \item Continuous batching (process requests without gaps)
\end{itemize}

\subsection{Era 4: The Commercial Deployment Era (2022--Present)}

ChatGPT's launch in November 2022 caused a step-change in the scale
of inference systems. Suddenly, systems that researchers had designed
for thousands of queries per day needed to handle \textit{millions}.

This era saw:
\begin{itemize}
  \item Purpose-built inference engines: \textbf{vLLM}, \textbf{TensorRT-LLM},
        \textbf{Triton Inference Server}, \textbf{TGI (Text Generation Inference)}
  \item New serving paradigms: \term{PagedAttention},
        \term{speculative decoding}, \term{disaggregated prefill}
  \item Hardware specialisation: TPUs, AWS Trainium/Inferentia,
        custom inference chips
  \item Multi-model serving, router architectures, mixture-of-experts
\end{itemize}

\begin{figure}[htbp]
\centering
\begin{tikzpicture}[xscale=1.35, yscale=0.85]

  % Timeline axis
  \draw[->, very thick, chaptercolor] (0,0) -- (11.5,0)
    node[right] {\textbf{Year}};
  \draw[->, very thick, chaptercolor] (0,0) -- (0,7.5)
    node[above] {\textbf{Scale}};

  % Y-axis labels
  \foreach \y/\label in {1/1M, 2/10M, 3/100M, 4/1B, 5/10B, 6/100B, 7/1T} {
    \draw[gray!30] (0,\y) -- (11.2,\y);
    \node[left, font=\tiny\color{gray}] at (-0.1,\y) {\label};
  }

  % X-axis ticks and labels
  \foreach \x/\year in {0.5/2012, 1.5/2014, 2.5/2016, 3.5/2018,
                         4.5/2019, 5.5/2020, 6.5/2021, 7.5/2022,
                         8.5/2023, 9.5/2024, 10.5/2025} {
    \draw[gray!40] (\x, -0.1) -- (\x, 0.1);
    \node[below, font=\tiny\color{gray}, rotate=45] at (\x,-0.15) {\year};
  }

  % Model parameter curve (approximate)
  \draw[chaptercolor, very thick, smooth]
    plot coordinates {
      (0.5, 1.5)   % AlexNet ~60M
      (1.5, 2.0)   % VGG ~140M
      (2.5, 2.5)   % ResNet ~25M (smaller but deeper)
      (3.5, 3.1)   % BERT 340M
      (4.5, 3.7)   % GPT-2 1.5B
      (5.5, 5.2)   % GPT-3 175B
      (6.5, 5.5)   % ~400B range
      (7.5, 5.8)   % ChatGPT era
      (8.5, 6.3)   % GPT-4 ~1.76T (estimated)
      (9.5, 6.5)
      (10.5,6.8)
    };

  % Era bands
  \fill[orange!8] (0,0) rectangle (3.0,7.2);
  \fill[yellow!8] (3.0,0) rectangle (5.0,7.2);
  \fill[green!8]  (5.0,0) rectangle (7.0,7.2);
  \fill[blue!8]   (7.0,0) rectangle (11.2,7.2);

  \node[font=\tiny\bfseries\color{orange!70}, rotate=90] at (1.5, 3.6)
    {Era 1: Research};
  \node[font=\tiny\bfseries\color{yellow!70!black}, rotate=90] at (4.0, 3.6)
    {Era 2: Pre-trained};
  \node[font=\tiny\bfseries\color{accentgreen!80}, rotate=90] at (6.0, 3.6)
    {Era 3: LLM};
  \node[font=\tiny\bfseries\color{chaptercolor}, rotate=90] at (9.1, 3.6)
    {Era 4: Commercial Scale};

  % Landmark model annotations
  \fill[orange!80] (0.5, 1.5) circle (3pt)
    node[above right, font=\tiny] {AlexNet};
  \fill[chaptercolor] (3.5, 3.1) circle (3pt)
    node[above right, font=\tiny] {BERT};
  \fill[chaptercolor] (5.5, 5.2) circle (3pt)
    node[above right, font=\tiny] {GPT-3};
  \fill[accentorange] (7.5, 5.8) circle (3pt)
    node[above right, font=\tiny] {ChatGPT};
  \fill[red!70] (8.5, 6.3) circle (3pt)
    node[above right, font=\tiny] {GPT-4 ($\sim$1.76T)};

  % Y-axis label (rotated)
  \node[rotate=90, font=\small\color{chaptercolor}] at (-1.2,3.6)
    {Approximate Parameters};

\end{tikzpicture}
\caption{Evolution of AI model scale over time, showing the four eras
of AI inference infrastructure. Note the logarithmic scale on the
y-axis --- each step represents a $10\times$ increase in parameters.}
\label{fig:evolution}
\end{figure}

%% -----------------------------------------------------------
\section{Modern Inference Systems: A Tour}
%% -----------------------------------------------------------

Let us briefly survey the major inference engines and systems
that power today's AI products.
We will study each of these in detail in later chapters.

\subsection{vLLM: PagedAttention and Continuous Batching}

vLLM, developed at UC Berkeley in 2023, is arguably the most
influential open-source LLM inference engine.
Its two key innovations are:

\begin{itemize}
  \item \term{PagedAttention}: A memory management technique that
        treats GPU memory like virtual memory in operating systems,
        reducing memory waste from 60--80\% down to under 4\%.
  \item \term{Continuous Batching}: Instead of waiting for a full
        batch before processing, vLLM dynamically adds and removes
        requests from the processing queue, dramatically improving GPU
        utilisation.
\end{itemize}

In benchmarks, vLLM achieves \textbf{2--4$\times$ higher throughput}
compared to naive serving approaches.

\subsection{TensorRT-LLM: NVIDIA's Production Engine}

TensorRT-LLM is NVIDIA's production inference engine.
It compiles model graphs into highly optimised machine code specific
to NVIDIA GPU architectures.

Key features:
\begin{itemize}
  \item Kernel fusion (combines multiple operations into one GPU kernel)
  \item INT8/FP8 quantisation with hardware-native support
  \item In-flight batching (similar to continuous batching)
  \item Custom CUDA kernels for attention, LayerNorm, etc.
\end{itemize}

TensorRT-LLM typically achieves \textbf{40--60\% higher throughput}
than PyTorch-based serving.

\subsection{Text Generation Inference (TGI)}

Hugging Face's TGI is widely used for open-source model deployment.
It supports a wide range of models and provides production-grade
features like tensor parallelism, quantisation, and Prometheus metrics.

\subsection{Triton Inference Server}

NVIDIA's Triton is a \textit{model-agnostic} serving framework that
supports TensorRT, PyTorch, TensorFlow, ONNX, and custom backends.
It is commonly used in enterprise deployments where multiple model
types need to be served from a unified platform.

\subsection{Production Cloud Services}

At the cloud provider level:
\begin{itemize}
  \item \textbf{AWS Bedrock}: Serves Claude, Llama, Mistral, and others
        on Inferentia2 and GPU instances
  \item \textbf{Google Vertex AI}: Serves Gemini on custom TPU infrastructure
  \item \textbf{Azure OpenAI}: Serves GPT-4 on custom NVIDIA clusters
  \item \textbf{Groq}: Uses custom LPU (Language Processing Unit) hardware
        designed specifically for LLM inference, achieving extremely
        high token generation speeds ($>$ 500 tokens/second)
\end{itemize}

%% -----------------------------------------------------------
\section{Real-World Applications and Use Cases}
%% -----------------------------------------------------------

Let us ground these concepts by looking at how real systems use
AI inference.

\subsection{Conversational AI (ChatGPT, Claude, Gemini)}

The canonical inference use case.
Key characteristics:
\begin{itemize}
  \item Interactive: users expect sub-second TTFT
  \item Variable length: inputs range from 10 to 100,000+ tokens
  \item Stateless per request (context is re-sent each time)
  \item Requires streaming output (word by word display)
  \item High concurrency: millions of simultaneous users
\end{itemize}

\subsection{Code Generation (GitHub Copilot, Cursor)}

Code completion requires:
\begin{itemize}
  \item Extremely low latency ($<$ 100ms) --- must feel instant as you type
  \item High throughput --- millions of suggestions per minute
  \item Short outputs (typically 1--10 lines)
  \item Context window efficiency (large code files as context)
\end{itemize}

GitHub Copilot reportedly handles over \textbf{1 million suggestions
per day} at sub-100ms latency.

\subsection{Document Processing (RAG, Summarisation)}

Retrieval-Augmented Generation (RAG) pipelines combine AI inference
with database retrieval:
\begin{enumerate}
  \item User query \textrightarrow{} embedding model inference (fast, $<$10ms)
  \item Embedding \textrightarrow{} vector database search
  \item Retrieved documents + query \textrightarrow{} LLM inference (slower, 1--10s)
\end{enumerate}

This requires \textit{two} types of inference: embedding models
(small, fast, batch-friendly) and generation models (large, sequential).

\subsection{Image Generation (Stable Diffusion, Midjourney, DALL-E)}

Diffusion models have unique inference characteristics:
\begin{itemize}
  \item Inference requires 20--50 \textit{denoising steps}
  \item Each step is a full forward pass through a large model
  \item Output is a 2D image, not text tokens
  \item Parallelism is easier (image patches can be processed in parallel)
  \item Users tolerate 5--30 second wait times
\end{itemize}

\subsection{Real-Time Applications (Voice, AR/VR)}

The frontier of inference engineering:
\begin{itemize}
  \item Voice assistants require end-to-end latency $<$ 300ms
        (speech-to-text + LLM + text-to-speech pipeline)
  \item AR/VR scene understanding needs 60+ FPS inference
  \item These require on-device inference (edge deployment)
\end{itemize}

%% -----------------------------------------------------------
\section{Your First Inference System: A Practical Example}
%% -----------------------------------------------------------

Let us make these concepts concrete with a simple but complete
Python example.
We will serve a language model and measure the key performance metrics.

\begin{lstlisting}[style=pythonstyle,
  caption={Simple LLM inference with performance measurement},
  label={lst:basic-inference}]
"""
Chapter 1: Basic LLM Inference Example
---------------------------------------
This script demonstrates the fundamentals of LLM inference:
  - Loading a model
  - Running a forward pass (inference)
  - Measuring latency and throughput
  - Understanding time-to-first-token vs total generation time

Run with: python basic_inference.py
Requires: pip install transformers torch
"""

import time
import torch
from transformers import AutoModelForCausalLM, AutoTokenizer

# -- Step 1: Load the model and tokeniser ------------------------------
# The tokeniser converts text to token IDs (numbers the model understands)
# The model contains all the trained parameters (weights)
#
# For this example we use a small model to avoid needing large GPU memory.
# In production, you would use GPT-4, LLaMA-3, Mistral, etc.

MODEL_NAME = "gpt2"  # Small model for illustration (~117M parameters)

print(f"Loading model: {MODEL_NAME}")
tokeniser = AutoTokenizer.from_pretrained(MODEL_NAME)
model = AutoModelForCausalLM.from_pretrained(
    MODEL_NAME,
    torch_dtype=torch.float16,  # Use 16-bit precision (half memory of float32)
    device_map="auto",          # Automatically place on GPU if available
)
model.eval()  # Set to inference mode (disables dropout, gradient tracking)

print(f"Model loaded. Parameters: {sum(p.numel() for p in model.parameters()):,}")
print(f"Device: {next(model.parameters()).device}")
print()

# -- Step 2: Define our inference function ----------------------------
def run_inference(prompt: str, max_new_tokens: int = 100) -> dict:
    """
    Runs a single inference request and returns the result with metrics.

    Returns a dict containing:
      - output_text:    the generated text
      - input_tokens:   number of tokens in the prompt
      - output_tokens:  number of tokens generated
      - ttft_ms:        time to first token (milliseconds)
      - total_ms:       total generation time (milliseconds)
      - tokens_per_sec: throughput (tokens/second)
    """

    # Step 2a: Tokenise the input
    # Convert raw text \textrightarrow{} token IDs \textrightarrow{} tensor on GPU
    inputs = tokeniser(prompt, return_tensors="pt").to(model.device)
    input_len = inputs["input_ids"].shape[1]  # number of input tokens

    # Step 2b: Record start time for TTFT measurement
    t_start = time.perf_counter()

    # Step 2c: Run inference
    # torch.no_grad() tells PyTorch NOT to track gradients
    # (we don't need them for inference, and they would waste memory)
    with torch.no_grad():
        output_ids = model.generate(
            **inputs,
            max_new_tokens=max_new_tokens,
            do_sample=False,     # Greedy decoding (deterministic)
            pad_token_id=tokeniser.eos_token_id,
        )

    t_end = time.perf_counter()

    # Step 2d: Compute metrics
    output_len = output_ids.shape[1] - input_len  # tokens we generated
    total_ms = (t_end - t_start) * 1000           # total time in ms
    tokens_per_sec = output_len / (t_end - t_start)

    # Step 2e: Decode the output back to text
    generated_tokens = output_ids[0, input_len:]  # only new tokens
    output_text = tokeniser.decode(generated_tokens, skip_special_tokens=True)

    return {
        "output_text":    output_text,
        "input_tokens":   input_len,
        "output_tokens":  output_len,
        "total_ms":       total_ms,
        "tokens_per_sec": tokens_per_sec,
    }


# -- Step 3: Run some example requests --------------------------------
prompts = [
    "Explain what artificial intelligence is in simple terms:",
    "The key difference between training and inference in machine learning is:",
    "A GPU is better than a CPU for deep learning because:",
]

for i, prompt in enumerate(prompts, 1):
    print(f"{'='*60}")
    print(f"Request {i}: {prompt[:60]}...")
    print()

    result = run_inference(prompt, max_new_tokens=80)

    print(f"  Input tokens:  {result['input_tokens']}")
    print(f"  Output tokens: {result['output_tokens']}")
    print(f"  Total time:    {result['total_ms']:.1f} ms")
    print(f"  Throughput:    {result['tokens_per_sec']:.1f} tokens/sec")
    print()
    print(f"  Generated text:")
    print(f"  {result['output_text'][:200]}")
    print()

# -- Step 4: Batch inference comparison -------------------------------
print(f"{'='*60}")
print("Comparing single vs batch inference:")
print()

# Run requests one-by-one
t_sequential_start = time.perf_counter()
for prompt in prompts:
    run_inference(prompt, max_new_tokens=50)
t_sequential = time.perf_counter() - t_sequential_start

print(f"  Sequential (3 requests): {t_sequential*1000:.0f} ms")
print()
print("  Note: In production, continuous batching processes multiple")
print("  requests simultaneously, dramatically improving throughput.")
print("  We will cover this in Chapter 4.")
\end{lstlisting}

Running this script on an NVIDIA A100 GPU with GPT-2 produces output
similar to:

\begin{lstlisting}[style=bashstyle,
  caption={Sample output from basic inference script}]
Loading model: gpt2
Model loaded. Parameters: 117,222,400
Device: cuda:0

============================================================
Request 1: Explain what artificial intelligence is in simpl...

  Input tokens:  12
  Output tokens: 80
  Total time:    245.3 ms
  Throughput:    326.1 tokens/sec

============================================================
Comparing single vs batch inference:
  Sequential (3 requests): 712 ms
\end{lstlisting}

\begin{notebox}[What to Observe]
Notice that 80 tokens took 245ms total. This is about 3ms per token.
On a larger model (e.g. LLaMA-3-70B), you would see:
\begin{itemize}
  \item Much higher memory usage (70B params $\times$ 2 bytes = 140GB)
  \item Slower throughput (perhaps 20--50 tokens/sec per request)
  \item Need for multi-GPU setup
\end{itemize}
These numbers motivate everything we will optimise in this book.
\end{notebox}

%% -----------------------------------------------------------
\section{Performance Measurement: The Metrics That Matter}
%% -----------------------------------------------------------

Before we end this chapter, we need to establish a shared vocabulary
for measuring inference system performance.
These metrics will appear throughout every chapter of this book.

\subsection{Latency Metrics}

\begin{defbox}[Time to First Token (TTFT)]
The time elapsed from when the request is received to when the
\textit{first output token} is generated.\\[4pt]
Formula: $\text{TTFT} = t_{\text{first\_token}} - t_{\text{request\_received}}$\\[4pt]
Typical targets: $<200$ms for interactive applications.
\end{defbox}

\begin{defbox}[Time Per Output Token (TPOT)]
The average time to generate each output token after the first one.
This is the inverse of the per-request generation speed.\\[4pt]
Formula: $\text{TPOT} = \frac{t_{\text{last\_token}} - t_{\text{first\_token}}}{N_{\text{output\_tokens}} - 1}$\\[4pt]
Typical targets: 20--50ms/token for single-stream, lower with batching.
\end{defbox}

\begin{defbox}[End-to-End Latency]
Total time from request submission to complete response received.\\[4pt]
Formula: $\text{E2E Latency} = \text{TTFT} + \text{TPOT} \times (N_{\text{output\_tokens}} - 1)$\\[4pt]
Used for SLA (Service Level Agreement) commitments.
\end{defbox}

\subsection{Throughput Metrics}

\begin{defbox}[Requests Per Second (RPS)]
How many complete inference requests the system can handle per second
across all concurrent users.\\[4pt]
This is the primary \textit{capacity} metric for inference systems.
\end{defbox}

\begin{defbox}[Output Tokens Per Second (TPS)]
Total output tokens generated per second across all requests.\\[4pt]
Formula: $\text{TPS} = \sum_{\text{requests}} \frac{N_{\text{output\_tokens}}}{t_{\text{generation}}}$\\[4pt]
This is the primary metric for comparing inference engine efficiency.
\end{defbox}

\subsection{Efficiency Metrics}

\begin{defbox}[GPU Utilisation]
Percentage of time the GPU's compute units are actively performing
calculations (as opposed to waiting for memory or network).\\[4pt]
Target: $>70\%$ in production systems. Below 50\% suggests the system
is bottlenecked by memory bandwidth or request scheduling.
\end{defbox}

\begin{defbox}[Cost Per Million Tokens (\$/MTok)]
The operational cost to generate one million output tokens.
Includes GPU compute, memory, networking, and power costs.\\[4pt]
This is the primary \textit{business} metric for inference operations.
Industry targets: \$0.10--\$5.00 per million tokens depending on
model size and infrastructure.
\end{defbox}

\begin{figure}[htbp]
\centering
\begin{tikzpicture}
\begin{axis}[
  width=0.85\textwidth, height=6cm,
  xlabel={Concurrent Requests},
  ylabel={Tokens per Second (System Total)},
  title={Throughput vs.\ Concurrency for Different Serving Strategies},
  legend pos=north west,
  grid=both,
  grid style={dotted, gray!40},
  xmin=0, xmax=32,
  ymin=0, ymax=2600,
  xtick={0,4,8,12,16,20,24,28,32},
  ytick={0,500,1000,1500,2000,2500},
  axis lines=left,
  tick align=outside,
  legend style={font=\small},
]

% Naive serving (one at a time)
\addplot[chaptercolor, thick, mark=square*, mark size=2pt]
  coordinates {
    (1,120) (2,138) (4,155) (8,162) (12,163) (16,164) (20,165) (32,164)
  };
\addlegendentry{Naive (sequential)};

% Static batching
\addplot[accentorange, thick, mark=triangle*, mark size=2pt]
  coordinates {
    (1,120) (2,230) (4,440) (8,820) (12,1100) (16,1280) (20,1350) (32,1380)
  };
\addlegendentry{Static batching};

% Continuous batching (vLLM-style)
\addplot[accentgreen, thick, mark=*, mark size=2pt]
  coordinates {
    (1,120) (2,240) (4,470) (8,900) (12,1400) (16,1800) (20,2100) (32,2450)
  };
\addlegendentry{Continuous batching};

\end{axis}
\end{tikzpicture}
\caption{Throughput scaling with concurrent requests for three serving
strategies. Continuous batching (vLLM-style) maintains near-linear
scaling to much higher concurrency levels. We will explain exactly
why in Chapter 4.}
\label{fig:throughput-vs-concurrency}
\end{figure}

%% -----------------------------------------------------------
\section{Production Engineering Notes}
%% -----------------------------------------------------------

Before we close this chapter, here are some production realities
that you will encounter as an inference engineer.
These are lessons learned from running real systems at scale.

\subsection{The Cold Start Problem}

When a model is first loaded onto a GPU (or after a GPU is restarted),
the first few requests are significantly slower because:
\begin{itemize}
  \item Model weights must be transferred from disk/CPU to GPU memory
  \item CUDA kernels must be compiled and cached
  \item Memory allocators warm up
\end{itemize}

Production systems use \term{warm-up requests} --- dummy requests sent
at startup to initialise the system before real traffic arrives.

\subsection{Memory Fragmentation}

Even with sophisticated memory management, GPU memory can become
fragmented over time --- unable to allocate large contiguous blocks
even when there is enough total free memory.
Production systems monitor memory fragmentation and periodically
restart workers to clean up.

\subsection{Long-Tail Latency}

The average latency might be excellent (200ms) but the P99 latency
(worst 1\% of requests) might be terrible (5 seconds).
This happens due to:
\begin{itemize}
  \item Garbage collection pauses
  \item Memory pressure causing evictions
  \item Network jitter
  \item Unlucky scheduling
\end{itemize}

Always monitor and design to P99, not just averages.

\subsection{Model Versioning and Rollouts}

Updating a model in production while maintaining availability requires
careful blue-green deployment or canary rollouts.
With large models (hundreds of GB), the update itself takes minutes ---
during which the new model must be loaded without dropping service.

%% -----------------------------------------------------------
\section{Common Mistakes in Inference Engineering}
%% -----------------------------------------------------------

\begin{warningbox}[Mistake 1: Confusing Training Metrics with Inference Metrics]
A model with excellent training loss does not necessarily produce
good inference performance. Training optimises model quality.
Inference engineering optimises system performance.
They are separate concerns.
\end{warningbox}

\begin{warningbox}[Mistake 2: Ignoring Memory Bandwidth]
Most inference engineers initially focus on compute (FLOPS).
In practice, LLM inference is almost always
\textit{memory-bandwidth-bound}, not compute-bound.
Optimising for the wrong bottleneck wastes engineering effort.
\end{warningbox}

\begin{warningbox}[Mistake 3: Benchmarking with Unrealistic Inputs]
Benchmarking with short prompts and single requests gives misleading
results. Always benchmark with realistic input distributions,
concurrency levels, and request patterns that match your actual
production workload.
\end{warningbox}

\begin{warningbox}[Mistake 4: Under-provisioning for Peak Load]
AI inference traffic is highly spiky --- product launches, viral
moments, and business hours can cause $10\times$ normal traffic in
minutes. Always provision for 2--3$\times$ expected peak, and have
clear shedding/queueing strategies for higher load.
\end{warningbox}

%% -----------------------------------------------------------
\section{Best Practices}
%% -----------------------------------------------------------

\begin{tipbox}[Best Practice 1: Measure Before Optimising]
Before optimising anything, measure your current system's bottleneck.
Use profiling tools (NVIDIA Nsight, PyTorch Profiler, vLLM's built-in
metrics) to determine whether you are compute-bound, memory-bound,
or I/O-bound. Then optimise the actual bottleneck.
\end{tipbox}

\begin{tipbox}[Best Practice 2: Use the Right Hardware]
Not all GPUs are equal for inference.
\begin{itemize}
  \item For memory-bandwidth-limited workloads (most LLMs): choose
        GPUs with highest HBM bandwidth (H100, A100)
  \item For compute-limited workloads (diffusion models, small models):
        choose GPUs with highest FLOPS
  \item For cost-sensitive deployments: consider A10G, L4 --- cheaper
        with lower throughput, acceptable for low-concurrency workloads
\end{itemize}
\end{tipbox}

\begin{tipbox}[Best Practice 3: Monitor the Full Stack]
Monitor metrics at every layer:
\begin{itemize}
  \item Application: TTFT, E2E latency, error rates
  \item Inference engine: batch sizes, queue depth, cache hit rate
  \item GPU: utilisation, memory usage, memory bandwidth
  \item Infrastructure: network latency, disk I/O
\end{itemize}
Slow P99 latency could be caused by a problem at \textit{any} layer.
\end{tipbox}

%% -----------------------------------------------------------
\section{Chapter Summary}
%% -----------------------------------------------------------

\begin{chapsummary}
\textbf{Key concepts from Chapter 1:}

\begin{itemize}
  \item \textbf{AI Inference} is the process of running a trained model
        on new data to produce outputs. The model parameters are fixed;
        only the input varies.

  \item \textbf{Training vs Inference} are fundamentally different:
        training changes parameters (slow, done once);
        inference applies fixed parameters (fast, done millions of times).

  \item \textbf{The inference request lifecycle} involves: network
        ingress, scheduling, tokenisation, prefill (parallel),
        decode (sequential), and streaming output.

  \item \textbf{The inference trilemma}: you must balance latency,
        throughput, and cost. Optimising all three simultaneously
        is the core challenge of the discipline.

  \item \textbf{Hardware matters}: GPUs dominate because neural networks
        are fundamentally massive parallel matrix operations.
        Memory bandwidth --- not raw compute --- is the primary
        bottleneck for most LLM inference.

  \item \textbf{Modern inference engines} (vLLM, TensorRT-LLM, TGI)
        provide production-grade features: continuous batching,
        KV caching, quantisation, and tensor parallelism.

  \item \textbf{Key metrics}: TTFT, TPOT, E2E latency, tokens per
        second, GPU utilisation, and \$/MTok.
\end{itemize}

\vspace{6pt}
In the next chapter, we will go deep on the transformer architecture ---
understanding exactly \textit{what} the GPU is computing during inference
and \textit{why} the model is structured the way it is.
\end{chapsummary}

%% -----------------------------------------------------------
\section{Exercises}
%% -----------------------------------------------------------

\begin{exercise}[Measuring Inference Metrics]
Run the code example from Listing~\ref{lst:basic-inference} with GPT-2.
Then modify it to:
\begin{enumerate}
  \item Measure the time for just the tokenisation step vs the model
        forward pass.
  \item Plot throughput (tokens/sec) as a function of output sequence
        length (try 10, 50, 100, 200, 500 tokens).
  \item Calculate the cost estimate: assume the machine costs
        \$3/hour to run. What is your \$/MTok?
\end{enumerate}
\end{exercise}

\begin{exercise}[Understanding the Memory Constraint]
GPT-3 has 175 billion parameters.
\begin{enumerate}
  \item If stored in FP32 (4 bytes per parameter), how many GB of
        memory does it require?
  \item In FP16 (2 bytes per parameter)?
  \item In INT8 (1 byte per parameter)?
  \item An H100 GPU has 80GB of VRAM. How many H100s would you
        need to hold GPT-3 in FP16?
  \item Research online: how many GPUs does OpenAI reportedly use
        to serve GPT-3/GPT-4? Does this match your calculation?
\end{enumerate}
\end{exercise}

\begin{exercise}[The Trilemma in Practice]
Consider three hypothetical AI serving systems, each optimising for
a different corner of the latency--throughput--cost trilemma:
\begin{enumerate}
  \item A real-time voice assistant (optimise for?)
  \item A batch document processing pipeline running nightly
        (optimise for?)
  \item A startup API offering cheap LLM access
        (optimise for?)
\end{enumerate}
For each, identify which metric is \textit{most} and \textit{least}
important, and describe one concrete engineering decision they would
make differently.
\end{exercise}

\begin{exercise}[Architecture Exploration]
Draw (on paper or digitally) the full request lifecycle diagram
from Figure~\ref{fig:request-lifecycle} from memory. Then:
\begin{enumerate}
  \item Identify which stage you think takes the longest. Why?
  \item Identify which stage is hardest to scale horizontally. Why?
  \item If you needed to reduce TTFT by 50\%, which stage would
        you focus on?
\end{enumerate}
\end{exercise}

%% -----------------------------------------------------------
\section{Further Reading}
%% -----------------------------------------------------------

\begin{furtherreading}
\textbf{Foundational Papers}
\begin{itemize}
  \item Vaswani et al.\ (2017). \textit{Attention Is All You Need}.
        The original transformer paper. Understanding this is
        prerequisites for Chapter 2.
  \item Kwon et al.\ (2023). \textit{Efficient Memory Management for
        Large Language Model Serving with PagedAttention}.
        The paper that introduced vLLM's core innovation.
  \item Pope et al.\ (2022). \textit{Efficiently Scaling Transformer
        Inference}. Google's analysis of LLM inference bottlenecks.
\end{itemize}

\textbf{Engineering Resources}
\begin{itemize}
  \item vLLM documentation: \url{https://docs.vllm.ai}
  \item NVIDIA TensorRT-LLM: \url{https://github.com/NVIDIA/TensorRT-LLM}
  \item Hugging Face TGI: \url{https://github.com/huggingface/text-generation-inference}
  \item Andrej Karpathy's ``LLM in 100 lines'' --- builds intuition
        for the transformer forward pass.
\end{itemize}

\textbf{Background Reading}
\begin{itemize}
  \item Dettmers et al.\ (2022). \textit{LLM.int8(): 8-bit Matrix
        Multiplication for Transformers at Scale}.
        Introduction to quantisation (Chapter 6 preview).
  \item Orca paper (Yu et al., 2022). \textit{Orca: A Distributed
        Serving System for Transformer-Based Language Generation Models}.
        Introduced continuous batching.
  \item Leviathan et al.\ (2023). \textit{Fast Inference from
        Transformers via Speculative Decoding}.
        A clever technique we will cover in Chapter 7.
\end{itemize}

\textbf{Online Resources}
\begin{itemize}
  \item \url{https://modal.com/blog} --- excellent practical articles
        on LLM serving
  \item \url{https://www.anyscale.com/blog} --- Ray Serve and LLM
        infrastructure deep dives
  \item MLCommons MLPerf Inference benchmarks ---
        \url{https://mlcommons.org/benchmarks/inference-datacenter/}
\end{itemize}
\end{furtherreading}

%% -- Chapter end ornament -----------------------------------
\vfill
\begin{center}
\begin{tikzpicture}
  \draw[chaptercolor!40, thick] (0,0) -- (4,0);
  \fill[chaptercolor!60] (2,0) circle (3pt);
  \draw[chaptercolor!40, thick] (4,0) -- (8,0);
\end{tikzpicture}

\vspace{6pt}
{\small\itshape\color{gray}
End of Chapter 1 $\cdot$
\textit{Next: Chapter 2 --- The Transformer Architecture: What the GPU Is Actually Computing}
}
\end{center}

\end{document}
